\documentclass[12pt]{article}

\usepackage[margin=1in]{geometry}
\usepackage[T1]{fontenc}
\usepackage[utf8]{inputenc}
\usepackage{lmodern}
\usepackage{hyperref}
\usepackage{longtable}
\usepackage{array}
\usepackage{booktabs}
\usepackage{enumitem}
\usepackage{xcolor}
\usepackage{fancyhdr}
\usepackage{verbatim}
\usepackage{titlesec}

\hypersetup{
  colorlinks=true,
  linkcolor=blue,
  urlcolor=blue,
  pdftitle={User Manual for Sanskrit Fragment Reconstruction Studio}
}

\setlist{itemsep=0.35em, topsep=0.35em}
\setlength{\parskip}{0.6em}
\setlength{\parindent}{0pt}
\setlength{\headheight}{15pt}

\pagestyle{fancy}
\fancyhf{}
\fancyhead[L]{Sanskrit Fragment Reconstruction Studio}
\fancyhead[R]{User Manual}
\fancyfoot[C]{\thepage}

\titleformat{\section}{\large\bfseries}{\thesection.}{0.5em}{}
\titleformat{\subsection}{\normalsize\bfseries}{\thesubsection}{0.5em}{}

\begin{document}

\pagenumbering{gobble}

\begin{titlepage}
\centering
{\Large Sanskrit Ciphers\par}
\vspace{1.5cm}
{\Huge\bfseries User Manual\par}
\vspace{0.6cm}
{\LARGE Sanskrit Fragment Reconstruction Studio\par}
\vspace{1.5cm}
{\large Professional User Guide for the Desktop Application\par}
\vfill

\begin{tabular}{>{\bfseries}p{4cm}p{9cm}}
Application & Sanskrit Fragment Reconstruction Studio \\
Project & Sanskrit Ciphers \\
Audience & Scholars, researchers, and non-technical users \\
Prepared From & Source code in the project's \texttt{src/} directory \\
Date & March 17, 2026 \\
\end{tabular}

\vfill
This manual describes the current desktop application implemented in the repository's
\texttt{src/web/web-canvas} module. Supporting machine learning tools are referenced
when they affect normal use of the application.
\end{titlepage}

\clearpage
\pagenumbering{roman}
\tableofcontents
\newpage
\pagenumbering{arabic}

\section{Introduction}

\subsection{Purpose of the Application}
Sanskrit Fragment Reconstruction Studio is a desktop application for organizing,
searching, and visually arranging manuscript fragment images. It is designed to help
users examine fragment images, place them on a reconstruction canvas, edit metadata,
and save reconstruction sessions for later review.

\subsection{Target Users}
This application is intended for:
\begin{itemize}
  \item scholars working with manuscript fragments,
  \item research assistants organizing image collections, and
  \item users who need a visual workspace without writing code.
\end{itemize}

\subsection{Overview of Features}
The current application provides the following user-facing features:
\begin{itemize}
  \item a home screen with project history and fragment search,
  \item a searchable fragment library,
  \item an interactive canvas for arranging fragments,
  \item project saving and reloading,
  \item fragment uploads,
  \item editable fragment metadata,
  \item bulk metadata editing for multiple fragments,
  \item filter tools for narrowing large fragment collections,
  \item optional segmented-image display for supported fragments, and
  \item edge-match suggestions generated from previously prepared AI matching data.
\end{itemize}

\subsection{AI-Assisted Features}
The application includes several AI-assisted features that help users work more efficiently with large fragment collections:
\begin{itemize}
  \item \textbf{Segmentation support:} some fragments can be shown in a segmented view, where the fragment is isolated from the background for easier visual comparison.
  \item \textbf{Metadata enrichment:} fragment records may already include AI-generated information such as line count, script type, edge-related flags, circle detection, and scale information.
  \item \textbf{Edge-match suggestions:} the app can display ranked candidate matches for a selected fragment based on previously computed similarity data.
\end{itemize}

These AI features support the user's decision-making process. They are best treated as helpful suggestions rather than final scholarly conclusions.

\section{System Requirements}

\subsection{Hardware Requirements}
\begin{itemize}
  \item A desktop or laptop computer
  \item Mouse or trackpad for drag-and-drop work
  \item Enough free disk space for fragment images, the local database, and saved projects
\end{itemize}

\subsection{Software Requirements}
\begin{itemize}
  \item Windows, macOS, or Linux, depending on the packaged build provided to the user
  \item Permission to install a desktop application on the target machine
\end{itemize}

\subsection{Data Requirements}
The packaged application is assumed to include the fragment data needed for normal use.
Internally, the app uses:
\begin{itemize}
  \item a local SQLite database of fragments and projects, and
  \item a local image folder used by that database.
\end{itemize}

For the intended end user, these data files are assumed to be bundled with the installer or packaged application.

\section{Installation Guide}

\subsection{Recommended Setup}
The intended user is expected to receive a packaged desktop build of the application.
In this setup, the fragment database and bundled image data already come with the installed app.

\subsection{Step-by-Step Installation}
\begin{enumerate}
  \item Locate the installer or packaged application provided by your project team.
  \item Run the installer, or open the packaged application according to your operating system.
  \item Follow the on-screen installation steps.
  \item When installation is complete, launch \textbf{Sanskrit Fragment Reconstruction Studio}.
\end{enumerate}

\subsection{What the Installer Is Expected to Include}
The packaged version is assumed to include:
\begin{itemize}
  \item the desktop application,
  \item the bundled fragment image set,
  \item the local database used for search and project management, and
  \item any precomputed AI-generated metadata shipped with the release.
\end{itemize}

\subsection{Uninstalling the Application}
\begin{enumerate}
  \item Close the application.
  \item Open your operating system's installed applications list.
  \item Select \textbf{Sanskrit Fragment Reconstruction Studio}.
  \item Choose the uninstall or remove option.
  \item Follow the on-screen prompts.
\end{enumerate}

If your installation stores additional user-created content, your operating system may ask whether you also want to remove application data.

\section{Getting Started}

\subsection{How to Launch the App}
After installation, launch the app from your desktop shortcut, applications menu,
or packaged application folder. When the app opens, you will first see the home page.

\subsection{First-Time Usage Walkthrough}
\begin{enumerate}
  \item Review the project history panel on the left side of the home page.
  \item If you are starting fresh, click \textbf{New Project}.
  \item If fragments already exist in the local database, use the search bar to search by fragment ID.
  \item Open a project or create a new one to enter the canvas workspace.
  \item Drag fragments from the left sidebar onto the canvas.
  \item Arrange fragments, edit metadata, and save your progress.
\end{enumerate}

\subsection{What Happens on First Use}
Depending on how your environment was prepared, you may see one of two common starting states:
\begin{itemize}
  \item \textbf{Preloaded environment:} the fragment library is already populated, search works immediately, and older projects may appear in the history list.
  \item \textbf{Fresh user environment:} the bundled fragment library is present, but your own project history begins empty until you create projects.
\end{itemize}

\section{Features \& Usage}

\subsection{Home Screen}
The home screen is the application's starting point.

\subsubsection{Create a New Project}
\begin{enumerate}
  \item Click \textbf{New Project}.
  \item The app creates a new reconstruction session.
  \item You are taken directly to the canvas workspace.
\end{enumerate}

\subsubsection{Open an Existing Project}
\begin{enumerate}
  \item Find the project in the \textbf{Session History} panel.
  \item Click the project name.
  \item The saved canvas arrangement opens in the workspace.
\end{enumerate}

\subsubsection{Rename or Delete a Project}
\begin{enumerate}
  \item Move the pointer over a saved project.
  \item Use the rename icon to change its name.
  \item Use the delete icon to remove it.
  \item Confirm deletion when prompted.
\end{enumerate}

\subsubsection{Search for a Fragment}
\begin{enumerate}
  \item Enter a fragment ID or search text in the main search bar.
  \item Wait for autocomplete suggestions, if available.
  \item Press \textbf{Enter}, click \textbf{Search}, or choose a suggestion.
  \item The app opens the canvas workspace and filters the fragment library to the search results.
\end{enumerate}

\subsection{Fragment Sidebar}
The sidebar on the left side of the canvas workspace shows the available fragment library.

\subsubsection{Browse and Search the Library}
\begin{itemize}
  \item Use the sidebar search field to search within the current library.
  \item Use the sort control to switch between A-to-Z and Z-to-A ordering.
  \item If the dataset is large, more fragments load as needed.
\end{itemize}

\subsubsection{Select Multiple Fragments}
\begin{itemize}
  \item Use \textbf{Select all} to select all visible fragments.
  \item Use the \textbf{Auto-pair} option to automatically include a fragment's reverse side when selecting, when such pairing exists.
  \item Clear the selection using the clear control in the selection bar.
\end{itemize}

\subsubsection{Add Fragments to the Canvas}
\begin{enumerate}
  \item Drag a fragment from the sidebar to the canvas, or
  \item select multiple fragments and click \textbf{Add to Canvas}.
  \item The selected fragments appear in the workspace.
\end{enumerate}

\subsubsection{Edit or Delete Fragments from the Library}
\begin{itemize}
  \item Click a fragment entry to open its metadata panel.
  \item Select multiple fragments and choose \textbf{Edit Metadata} for bulk changes.
  \item Select multiple fragments and choose \textbf{Delete} to permanently remove them from the database.
\end{itemize}

\subsection{Canvas Workspace}
The canvas is the main visual reconstruction area.

\subsubsection{Basic Placement and Selection}
\begin{itemize}
  \item Drag fragments onto the canvas to place them.
  \item Click a fragment to select it.
  \item Use \textbf{Shift} or \textbf{Ctrl/Cmd} with click to select multiple fragments.
  \item Click an empty area to clear the current selection.
\end{itemize}

\subsubsection{Move, Resize, and Rotate}
\begin{itemize}
  \item Drag a fragment to move it.
  \item Use the transform handles to resize it.
  \item Use the rotation handle to rotate it.
  \item If grouped fragments are selected together, they move together.
\end{itemize}

\subsubsection{Toolbar Actions}
The top toolbar provides the main workspace controls.

\begin{longtable}{>{\raggedright\arraybackslash}p{3cm} p{10.5cm}}
\toprule
\textbf{Control} & \textbf{Purpose} \\
\midrule
Home & Return to the home screen \\
Lock / Unlock & Prevent or allow editing of selected fragments \\
Front / Back & Change visual stacking order \\
Group / Ungroup & Move selected fragments as a set or separate them again \\
Mirror & Flip one selected fragment horizontally \\
Segmented / Original & Switch between segmented and original image display for one supported fragment \\
Edit Metadata & Edit metadata for two or more selected fragments at once \\
Delete & Remove selected fragments from the canvas only \\
Grid & Show or hide the measurement grid \\
Filters & Open or close the filter panel \\
Clear Canvas & Remove all fragments from the canvas after confirmation \\
Save & Save the current canvas state immediately \\
Upload & Add new image files to the fragment library \\
\bottomrule
\end{longtable}

\subsubsection{Zoom and Navigation}
\begin{itemize}
  \item Use the mouse wheel or trackpad gesture to zoom.
  \item Drag on empty canvas space to pan.
  \item Use the \textbf{Recenter} button to center the view on the current fragments.
  \item Use the zoom controls in the lower-right corner to zoom in, zoom out, or reset zoom.
\end{itemize}

\subsubsection{Grid Display}
When the grid is turned on:
\begin{itemize}
  \item small squares represent 1 cm by 1 cm,
  \item bold grid intervals represent 5 cm by 5 cm, and
  \item the app briefly shows a scale reminder in the lower-left corner.
\end{itemize}

\subsubsection{Edge Match Suggestions}
When exactly one canvas fragment is selected, an \textbf{Edge Match?} button may appear.

To use it:
\begin{enumerate}
  \item Select one fragment already on the canvas.
  \item Click \textbf{Edge Match?}
  \item Review the ranked suggestions in the sidebar.
  \item Click \textbf{Place} beside a suggested match to place it on the canvas.
\end{enumerate}

This feature uses previously generated AI matching results stored in the packaged data. It helps users review likely neighboring fragments more quickly.

\subsection{Metadata Editing}

\subsubsection{Open Fragment Metadata}
You can open the metadata dialog in two ways:
\begin{itemize}
  \item click a fragment entry in the sidebar, or
  \item double-click a fragment already placed on the canvas.
\end{itemize}

\subsubsection{Fields You Can Edit}
The metadata dialog allows editing of user-facing fields such as:
\begin{itemize}
  \item line count,
  \item script type,
  \item transcription,
  \item notes,
  \item edge-related flags,
  \item circle detection flag,
  \item scale values,
  \item and any custom metadata fields that have been added by your team.
\end{itemize}

Some of these values may already be filled in by AI-assisted preprocessing. Users can review, correct, or replace them when needed.

\subsubsection{How to Edit a Field}
\begin{enumerate}
  \item Open the metadata dialog.
  \item Click the value or edit icon beside the field.
  \item Enter or choose the new value.
  \item Save the change.
\end{enumerate}

\subsubsection{Bulk Metadata Editing}
To update multiple fragments at once:
\begin{enumerate}
  \item select multiple fragments from the sidebar or canvas,
  \item click \textbf{Edit Metadata},
  \item check only the fields you want to update,
  \item enter the new values, and
  \item click \textbf{Apply}.
\end{enumerate}

\subsection{Filters}
The filter panel on the right side of the canvas workspace helps narrow the fragment library.

\subsubsection{Available Built-In Filters}
The current app supports filters for:
\begin{itemize}
  \item minimum and maximum line count,
  \item script type,
  \item edge-piece status,
  \item top edge,
  \item bottom edge,
  \item left edge,
  \item right edge,
  \item circle presence, and
  \item search text.
\end{itemize}

\subsubsection{Custom Filters}
Users can also create custom filters from the filter panel.

To create one:
\begin{enumerate}
  \item open the filter panel,
  \item click \textbf{Add Custom Filter},
  \item enter a filter name,
  \item choose either \textbf{Multiselect} or \textbf{Text match},
  \item add options if needed, and
  \item click \textbf{Create Filter}.
\end{enumerate}

Custom filters can later be edited or deleted from the same panel.

\subsubsection{Apply or Reset Filters}
\begin{itemize}
  \item Click \textbf{Apply Filters} to refresh the library using the current settings.
  \item Click \textbf{Reset Filters} to return to the default view.
\end{itemize}

\subsection{Uploading New Fragments}
The app can import image files directly into the local fragment library.

\begin{enumerate}
  \item Click \textbf{Upload} in the top toolbar.
  \item Click \textbf{Select Images}.
  \item Choose one or more JPG, JPEG, or PNG files.
  \item Review the selected file list.
  \item Click \textbf{Upload}.
  \item Wait for the result messages.
\end{enumerate}

Uploaded fragments appear in the sidebar after the library refreshes. Their metadata can be edited later. Newly uploaded files may not immediately have the same AI-generated metadata as the fragments that were bundled with the application.

\section{User Interface Overview}

\subsection{Main Screens}

\subsubsection{Home Page}
The home page includes:
\begin{itemize}
  \item a collapsible project history panel on the left,
  \item a central search bar with autocomplete,
  \item a \textbf{New Project} button,
  \item a count of active projects, and
  \item a count of available fragments.
\end{itemize}

\subsubsection{Canvas Page}
The canvas page includes:
\begin{itemize}
  \item a top toolbar for workspace actions,
  \item a left fragment sidebar,
  \item a central canvas area,
  \item a right filter panel, and
  \item pop-up dialogs for metadata and uploads.
\end{itemize}

\subsection{Navigation}
\begin{itemize}
  \item Use \textbf{Home} in the toolbar to return to the home page.
  \item Use the sidebar toggle buttons to hide or show the fragment and filter panels.
  \item Use search from the home page to jump directly into a filtered canvas session.
\end{itemize}

\subsection{Status Indicators}
The application shows several useful status indicators:
\begin{itemize}
  \item project save state: \textbf{Saved}, \textbf{Saving...}, or \textbf{Unsaved},
  \item active filter badge on the filter button,
  \item loading indicators while fragments are being fetched,
  \item upload results after image import, and
  \item search-result messages when search mode is active.
\end{itemize}

\section{Example Workflows}

\subsection{Workflow 1: Start a New Reconstruction}
\begin{enumerate}
  \item Launch the app.
  \item Click \textbf{New Project}.
  \item Use the fragment sidebar to find relevant fragments.
  \item Drag fragments onto the canvas.
  \item Move, rotate, and resize them until you reach a useful arrangement.
  \item Click \textbf{Save}.
\end{enumerate}

\subsection{Workflow 2: Search for a Known Fragment ID}
\begin{enumerate}
  \item Open the home page.
  \item Enter the fragment ID in the search bar.
  \item Choose the matching autocomplete result, or click \textbf{Search}.
  \item Review the filtered results in the canvas workspace.
  \item Drag the fragment onto the canvas and inspect its metadata if needed.
\end{enumerate}

\subsection{Workflow 3: Add New Fragment Images}
\begin{enumerate}
  \item Open any canvas project.
  \item Click \textbf{Upload}.
  \item Select one or more image files.
  \item Upload them.
  \item Wait for the sidebar to refresh.
  \item Open the metadata dialog and add useful descriptive information.
\end{enumerate}

\subsection{Workflow 6: Use AI Assistance During Review}
\begin{enumerate}
  \item Search for or place a fragment on the canvas.
  \item Check whether the fragment already includes AI-generated metadata such as line count, script type, or edge information.
  \item If available, switch to segmented view to reduce background distraction.
  \item Use \textbf{Edge Match?} to review suggested neighboring fragments.
  \item Accept, reject, or manually refine the arrangement based on your own judgment.
\end{enumerate}

\subsection{Workflow 4: Review Possible Matches}
\begin{enumerate}
  \item Place a fragment on the canvas.
  \item Select that single fragment.
  \item Click \textbf{Edge Match?}
  \item Review the ranked suggestions in the sidebar.
  \item Click \textbf{Place} on a promising candidate.
  \item Compare the placement visually and continue refining the arrangement.
\end{enumerate}

\subsection{Workflow 5: Standardize Metadata Across Many Fragments}
\begin{enumerate}
  \item Select several fragments from the sidebar.
  \item Click \textbf{Edit Metadata}.
  \item Enable only the fields you want to change.
  \item Enter the shared value.
  \item Apply the update.
\end{enumerate}

\section{Troubleshooting}

\subsection{The App Opens but No Fragments Appear}
Possible causes:
\begin{itemize}
  \item the installed package was incomplete,
  \item the bundled data was not installed correctly, or
  \item no fragments match the current filters.
\end{itemize}

Try this:
\begin{itemize}
  \item reset the filters,
  \item clear any search text,
  \item restart the app, or
  \item contact your project team to confirm that the packaged data was installed correctly.
\end{itemize}

\subsection{Search Returns No Results}
Possible causes:
\begin{itemize}
  \item the fragment ID does not exist in the database,
  \item the spelling is different from the stored record, or
  \item the current filters exclude the fragment you are trying to find.
\end{itemize}

\subsection{Uploaded Images Do Not Show Up Immediately}
\begin{itemize}
  \item Wait a moment for the upload result and refresh cycle.
  \item Check whether any file failed during upload.
  \item Reopen the upload dialog and try again if necessary.
\end{itemize}

\subsection{The Segmented/Original Toggle Does Not Appear}
This is normal when the fragment does not have segmentation data available. The toggle is shown only for fragments that already contain segmentation information.

\subsection{Edge Match Suggestions Do Not Appear}
This usually means the selected fragment does not have packaged edge-match data available, or the installed data set does not include that feature for the current fragment.

\subsection{A Project Was Not Saved}
The app includes auto-save, but users should still click \textbf{Save} before closing an important session. If the status indicator shows \textbf{Unsaved}, wait or save manually before exiting.

\subsection{A Metadata Value Is Rejected}
Some metadata fields are validated. For example, number-based fields and scale values must use valid numeric input, and script type must match one of the supported options.

\section{Limitations \& Assumptions}

\begin{itemize}
  \item This manual describes the desktop application, not the full machine learning training pipeline.
  \item The packaged release is assumed to include the data needed for normal use.
  \item Edge matching, segmentation display, and some metadata depend on AI-generated preprocessing performed before the app is delivered to the user.
  \item Project notes exist in the database schema, but they are not currently exposed as a normal end-user feature in the interface.
  \item The desktop app is intended for local use on one machine. Collaboration, cloud sync, and multi-user editing are not described by the present codebase.
  \item Uploaded images are added to the local environment only. The app does not provide a built-in remote backup service.
  \item AI-generated metadata and match suggestions are assistive features. They may be helpful, but they should still be reviewed by the user.
  \item Bulk deletion from the sidebar permanently removes fragments from the database and should be used carefully.
\end{itemize}

\section{FAQ}

\subsection{Do I need programming knowledge to use the app?}
No. The application is designed as a graphical desktop tool. Most tasks are completed with buttons, search fields, and drag-and-drop actions.

\subsection{Can I use the app if I only have image files?}
Yes. You can upload additional image files, but bundled fragments may include richer AI-generated metadata than newly uploaded files.

\subsection{Does the app save my work automatically?}
Yes. The app includes auto-save, but using the \textbf{Save} button is still recommended before closing a project.

\subsection{Can I edit metadata after uploading fragments?}
Yes. Open a fragment's metadata dialog to update line count, script type, transcription, notes, edge flags, scale values, and custom fields.

\subsection{Can I edit several fragments at once?}
Yes. Select multiple fragments, then use \textbf{Edit Metadata}.

\subsection{Why are some fragments shown with a segmented view and others are not?}
Segmented display is available only when segmented data was prepared for that fragment in the bundled application data.

\subsection{Can the app calculate matches by itself while I am using it?}
The app displays packaged edge-match suggestions that were generated ahead of time. In normal end-user use, it acts as a viewer for those AI results rather than generating new ones live.

\subsection{What is the safest way to close the app?}
Check that the save indicator shows \textbf{Saved}, or click \textbf{Save} manually before closing the window.

\end{document}
